\documentclass[a4paper]{article}
\usepackage[pdftex]{graphicx}
\usepackage[utf8]{inputenc}
\usepackage{enumerate}
\usepackage{icomma}
\usepackage{href-ul}
\hypersetup{
	colorlinks=true,
	linkcolor=blue,
	urlcolor=blue}
\usepackage{siunitx}
\sisetup{locale=DE} 
\usepackage{amssymb}
\usepackage{geometry}
\geometry{a4paper, top=15mm, left=15mm, right=15mm, bottom=15mm,
	headsep=10mm, footskip=12mm}

\begin{document}
{\bf Arbeit, Energie und Leistung }\\
\begin{enumerate}[1.]
%Aufgabe 1
{\bf \item Aufgabe}
\begin{enumerate}[a)]
\item Beschreibe kurz den Unterschied zwischen Arbeit und Leistung.
\item Begründe kurz folgende Aussage: ,,Wenn zwei die gleiche Arbeit verrichten, dann müssen sie noch lange nicht das Gleiche leisten".
\item Wie kann man feststellen, ob eine Person A (im physikalischen Sinn)
 mehr leistet als eine Person B ?
\end{enumerate}

{\item \bf Ein PKW besitze die Leistung $P=\SI{150}{\kilo\watt}$. - Welche Energie kann dieser Motor in einer Stunde abgeben?}

{\bf \item Aufgabe }
\begin{enumerate}[a)]
\item Erläutere folgende Aussagen:
\begin{itemize}
\item ,,Der Wirkungsgrad einer realen Maschine ist immer kleiner als 1 ."
\item ,,Ein Auto besitzt einen Wirkungsgrad von $\eta = 40 \%$."
\end{itemize}
\item Ein Kran muss für das Heben einer Kreissäge ($m=\SI{90}{\kilogram}$) auf $\SI{25}{\meter}$ Höhe eine Energie von etwa $\SI{25}{\kilo\joule}$ aufwenden. \\
Welchen Wirkungsgrad besitzt dieser Kran?
\item Die Säge hängt (üblicherweise) über das Wochenende am Kran. - Wieviel Arbeit verrichtet der Kran in dieser Zeit? (Kurze Begründung).  
\end{enumerate}



{\bf \item Auf welche Höhe wurde eine Palette mit Getränkekisten mit der Gewichtskraft von $F_G=\SI{1500}{\newton}$
gehoben,} wenn an ihr eine Arbeit von $W=\SI{12}{\kilo\joule}$ verrichtet wurde?

{\bf \item In einer ,,Sonnenfarm" sind Solarmodule so angebracht,} dass sie sich im Verlaufe eines Tages stets senkrecht zu den Sonnenstrahlen ausrichten. In der Sahara gelangt so 12 Stunden lang eine durchschnittliche Strahlungsleistung von $ \SI{1}{\kilo\watt\per\meter^2} $ auf die Module.

\begin{enumerate}
\item Wie viel Energie kann in dieser Zeit von einem Quadratmeter Mo\-dulfläche gewonnen werden? 
\item Wie groß ist die gewonnene Energie, wenn der Wirkungsgrad nur \\
$\eta = 17,5\%$ beträgt?
\item Wie groß müsste die Fläche einer solchen Sonnenfarm sein, wenn \\
$\SI{6e10}{\joule}$ an täglicher Energie gewonnen werden soll? 
\end{enumerate}


{\bf \item Eine Motorpumpe mit $P=\SI{2,0}{\kilo\watt}$ Leistung fördert in $t= \SI{4,0}{\minute}$ 5000 Liter Wasser aus einem Brunnenschacht}
\begin{enumerate}[a)]
\item Welche Art von mechanischer Arbeit verrichtet die Pumpe?
\item Wie tief ist der Schacht?
\end{enumerate}
\end{enumerate} 

\fbox{
	\begin{minipage}{0.5\textwidth}
		Zur Lösung bitte \href{https://www.okuyakl.de/phys/p8arbL401/ll401.pdf}{hier klicken} oder den QR-Code scannen.\\
	Weitere Arbeitsblätter gibt es unter 
	
	\href{https://www.okuyakl.de}{www.okuyakl.de}
	\end{minipage}
	\hfill
	\begin{minipage}{0.4\textwidth}
		\includegraphics[width=1.5 cm]{../../viecher/zwe03}
		\includegraphics[width=3 cm]{qr401}
		\includegraphics[width=2 cm]{../../viecher/afanticon1}
		
	\end{minipage}}

\end{document}
